% Requires \usepackage{amsmath,graphicx}.
% Adjust the graphics path if this file is moved into the paper repository.

The two certificates can also be viewed as radial subsets of configuration
space.  Fix the initial configuration $q_0$, clearance row $i$, and control
speed $\bar u$.  For each direction $u$ with $\lVert u\rVert_2=\bar u$, let
$\tau_{\mathrm{lip}}(q_0,u)$ and $\tau_{\mathrm{disp}}(q_0,u)$ denote the
times at which the corresponding displacement bounds first consume the
available clearance.  The regions plotted below are
\begin{align}
  \mathcal C_{\mathrm{lip}}(q_0)
  &= \bigl\{q_0+t u:
      \lVert u\rVert_2=\bar u,\ 0\le t\le
      \tau_{\mathrm{lip}}(q_0,u)\bigr\}, \\
  \mathcal C_{\mathrm{disp}}(q_0)
  &= \bigl\{q_0+t u:
      \lVert u\rVert_2=\bar u,\ 0\le t\le
      \tau_{\mathrm{disp}}(q_0,u)\bigr\}.
  \label{eq:radial-certificate-regions}
\end{align}
Thus, membership certifies the radial segment from $q_0$ to that configuration;
it does not assert that an arbitrary curve contained in the plotted set is
certified by a single held control.

\begin{figure}[t]
  \centering
  \includegraphics[width=\linewidth]{%
    results/certificate_horizon_explainer/certificate_horizons.pdf}
  \caption{Geometric comparison of the jointwise and motion-dependent
  certificates for a two-link planar arm.  (a)~A spherical point on the arm has
  initial clearance $h_i(q_0)$ and moves under the held control $u^\star$;
  ghosted poses show the resulting motion.  (b)~Evaluating each certificate over
  all control directions gives the weighted-$\ell_1$ region
  $\mathcal C_{\mathrm{lip}}(q_0)$ and the larger motion-dependent region
  $\mathcal C_{\mathrm{disp}}(q_0)$.  The shaded set
  $\mathcal C_{\mathrm{obs}}$ is the exact configuration obstacle for the
  spherical clearance row.  The highlighted ray is the particular trajectory
  $q(t)=q_0+t u^\star$.  (c)~Along that ray, the jointwise bound consumes the
  available clearance at $\tau_{\mathrm{lip}}$, whereas the integrated
  motion-dependent bound crosses it later at $\tau_{\mathrm{disp}}$.}
  \label{fig:certificate-regions}
\end{figure}

Figure~\ref{fig:certificate-regions} makes the source of the improvement
explicit.  The jointwise certificate charges the maximum displacement induced
by every joint independently, producing a weighted-$\ell_1$ diamond.  The
motion-dependent certificate retains the velocity induced by the particular
joint combination and bounds its growth along the held-control trajectory.  In
the illustrated direction,
\begin{equation}
  L_d D_{\mathrm{disp}}(t) \le D_{\mathrm{lip}}(t),
  \qquad
  \tau_{\mathrm{disp}} \ge \tau_{\mathrm{lip}},
\end{equation}
so the same state, control, and initial clearance certify a longer continuous
trajectory segment without an intermediate collision evaluation.
